\chapter{Происхождение центрального следового веса синглета и триплета}
\label{ch:v8-baryon-c0-singlet-triplet-central-trace-weight-parent-origin-gate}

Предыдущий гейт локализовал всю относительную свободу в центральном весе
\begin{equation}
 \rho_p=pP_{\mathbf1}+\frac{1-p}{3}P_{\mathbf3},
 \qquad 0<p<1,
 \qquad r(p)=\frac{1-p}{3p}.
 \label{eq:v8-baryon-c0-central-weight-parent-origin-family}
\end{equation}
Проверим, является ли какая-либо точка этого интервала следствием уже
построенного родительского действия.

\section{Условный счётный след не является происхождением}

Обычный нормированный след четырёхмерного представления даёт
\begin{equation}
 \rho_{\rm count}=\frac14I_4,
 \qquad p_{\rm count}=\frac14,
 \qquad r_{\rm count}=1.
 \label{eq:v8-baryon-c0-ambient-counting-trace-candidate}
\end{equation}
Однако на чётной алгебре $\mathbb C\oplus M_3(\mathbb C)$ любой допустимый
вес получается центральным переопределением того же следа:
\begin{equation}
 Z_p=4pP_{\mathbf1}+\frac{4(1-p)}{3}P_{\mathbf3},
 \qquad \rho_p=Z_p\rho_{\rm count},
 \qquad [Z_p,SO(3)]=[Z_p,\Gamma_4]=0.
 \label{eq:v8-baryon-c0-central-density-reweighting}
\end{equation}
В частности, равный полный вес двух неприводимых секторов даёт столь же
положительный и симметричный свидетель
\begin{equation}
 \rho_{\rm sector}=\frac12P_{\mathbf1}+\frac16P_{\mathbf3},
 \qquad p_{\rm sector}=\frac12,
 \qquad r_{\rm sector}=\frac13.
 \label{eq:v8-baryon-c0-central-weight-sector-witness}
\end{equation}
Следовательно, $p=1/4$ следует не из самого представления, а из
дополнительного требования одинаковой микроскопической следовой плотности
на четырёх стрелках. Текущий родитель этого требования не содержит: новая
семейная тройка, её состояние и endpoint-блок ранее получили реестр
\begin{equation}
 N_{\rm extension}=0/3.
 \label{eq:v8-baryon-c0-central-weight-unoriginated-extension-ledger}
\end{equation}

\section{Grading-суперслед и тепловой родитель}

Нормированный суперслед не исправляет ситуацию. При
\begin{equation}
 \Gamma_4=\operatorname{diag}(-1,+1,+1,+1),
 \qquad \frac{\Gamma_4}{\Tr\Gamma_4}
 =\frac12\operatorname{diag}(-1,+1,+1,+1)
 \label{eq:v8-baryon-c0-normalized-supertrace-density}
\end{equation}
первое собственное значение равно $-1/2$, поэтому это не положительное
состояние и не допустимый шумовой вес.

Наиболее общий $SO(3)$-инвариантный гамильтониан на носителе
$\mathbf1\oplus\mathbf3$ имеет вид
\begin{equation}
 h_{\rm cen}=\varepsilon_1P_{\mathbf1}
 +\varepsilon_3P_{\mathbf3},
 \qquad \Delta=\varepsilon_3-\varepsilon_1.
 \label{eq:v8-baryon-c0-central-two-level-hamiltonian}
\end{equation}
Его состояние Гиббса задаёт
\begin{equation}
 p(\beta\Delta)=\frac{1}{1+3e^{-\beta\Delta}},
 \qquad r(\beta\Delta)=e^{-\beta\Delta}.
 \label{eq:v8-baryon-c0-central-gibbs-weight-and-rate}
\end{equation}
Но это параметризация, а не селектор: для каждого $p\in(0,1)$ существует
\begin{equation}
 \beta\Delta(p)=\log\frac{3p}{1-p},
 \label{eq:v8-baryon-c0-central-gap-realizes-any-weight}
\end{equation}
и потому тепловой класс реализует весь исходный симплекс. Два прежних
свидетеля соответствуют разным допустимым щелям:
\begin{equation}
 p=\frac14\Longleftrightarrow\beta\Delta=0,
 \qquad
 p=\frac12\Longleftrightarrow\beta\Delta=\log3.
 \label{eq:v8-baryon-c0-central-gibbs-gap-witnesses}
\end{equation}
KMS связывает прямую и обратную скорости после выбора $h_{\rm cen}$, но не
выводит саму щель $\Delta$.

\section{Энтропийный near miss}

Энтропия центрального семейства равна
\begin{equation}
 S(p)=-p\log p-(1-p)\log\frac{1-p}{3}.
 \label{eq:v8-baryon-c0-central-weight-entropy}
\end{equation}
Её производные имеют точный вид
\begin{equation}
 S'(p)=\log\frac{1-p}{3p},
 \qquad
 S''(p)=-\frac{1}{p(1-p)}<0,
 \label{eq:v8-baryon-c0-central-weight-entropy-derivatives}
\end{equation}
так что максимум действительно единственен:
\begin{equation}
 \operatorname*{argmax}_{0<p<1}S(p)=\left\{\frac14\right\}.
 \label{eq:v8-baryon-c0-central-weight-maximum-entropy-candidate}
\end{equation}
Но функционал максимизации энтропии для connector-скоростей не является
членом текущего действия. Он повторяет условный счётный выбор, не объясняя
его происхождения.

Итоговый аудит источников равен
\begin{equation}
 N_{\rm parent\ origin}=0/6:
 \quad\{S_{\rm old},\Tr_4,\Str_{\Gamma},\text{KMS/Gibbs},
 S_{\max},\text{endpoint incidence}\}.
 \label{eq:v8-baryon-c0-central-weight-parent-origin-ledger}
\end{equation}

\begin{theorem}[Тепловая параметризация без происхождения центрального веса]
Каждый положительный центральный вес $p\in(0,1)$ представим единственным
безразмерным разрывом $\beta\Delta=\log(3p/(1-p))$. Поэтому KMS- или
Gibbs-условие не выбирает $p$ без независимого происхождения
singlet--triplet щели. Обычный след и максимум энтропии условно дают
$p=1/4$, но соответствующий принцип отсутствует в текущем родителе.
\end{theorem}

\begin{nogocriterion}[Нельзя считать размерность сектора динамическим весом]
Кратности $1$ и $3$ превращаются в веса $1/4$ и $3/4$ только после выбора
одинаковой следовой плотности на каждом состоянии. Центральное
переопределение $Z_p$ сохраняет семейную симметрию, grading,
положительность и нормировку, поэтому одна размерность представления не
является parent-origin.
\end{nogocriterion}

\begin{protocol}\begin{enumerate}
\item центральный симплекс: \textbf{$0<p<1$};
\item условный обычный след: \textbf{$p=1/4$, $r=1$};
\item центральное переопределение сохраняет симметрии: \textbf{да};
\item нормированный grading-суперслед положителен: \textbf{нет};
\item Gibbs-класс реализует все $p$: \textbf{да};
\item минимальный тепловой параметр: \textbf{$\beta\Delta$};
\item максимум энтропии условно выбирает: \textbf{$p=1/4$};
\item энтропийный принцип выведен: \textbf{нет};
\item parent-origin источников: \textbf{$0/6$};
\item физическая одиночная карта $c_0$: \textbf{не выведена};
\item следующий гейт: \textbf{минимальные гамильтоновы данные центрального веса}.
\end{enumerate}\end{protocol}

Машинный сертификат сохранён в
\path{s2t_v8_baryon_c0_singlet_triplet_central_trace_weight_parent_origin_gate_results.json}.